\documentclass[conference]{IEEEtran}
\IEEEoverridecommandlockouts

\usepackage{cite}
\usepackage{amsmath,amssymb,amsfonts}
\usepackage{algorithmic}
\usepackage{algorithm}
\usepackage{graphicx}
\usepackage{textcomp}
\usepackage{xcolor}
\usepackage{booktabs}
\usepackage{multirow}
\usepackage{url}

\def\BibTeX{{\rm B\kern-.05em{\sc i\kern-.025em b}\kern-.08em
    T\kern-.1667em\lower.7ex\hbox{E}\kern-.125emX}}

\begin{document}

\title{AdaDA-TransUNet: Adaptive Dual Attention Transformer UNet for Hardware Efficient Medical Image Segmentation}

\author{
\IEEEauthorblockN{Author Name}
\IEEEauthorblockA{\textit{Department of Computer Science} \\
\textit{University Name}\\
City, Country \\
author@university.edu}
}

\maketitle

\begin{abstract}
Medical image segmentation is a clinically critical task demanding both high spatial precision and computational efficiency for practical deployment. Transformer-based architectures, particularly those enhanced with dual attention mechanisms combining position attention and channel attention, have achieved state-of-the-art performance but suffer from quadratic complexity in both spatial and channel dimensions, rendering them inefficient on resource-constrained hardware. In this paper, we propose AdaDA-TransUNet, a hardware-aware adaptive dual attention Transformer UNet that addresses these limitations through four complementary innovations applied directly to the dual attention block: (1) a low-rank windowed position attention module that reduces spatial attention complexity from quadratic to sub-linear with respect to spatial resolution; (2) a grouped channel attention module that reduces channel attention complexity by a factor equal to the number of groups; (3) a differentiable adaptive feature gate that learns a per-channel blending ratio between position and channel attention, replacing the fixed compression ratio; and (4) a hardware-aware configuration module that automatically selects window size, rank, and group count based on available device memory at inference time. Experiments on three benchmark medical image segmentation datasets spanning CT, endoscopy, and dermatology modalities demonstrate that AdaDA-TransUNet achieves competitive accuracy compared to DA-TransUNet while substantially reducing training VRAM and enabling multi-GPU data-parallel training where the baseline fails. The proposed AdaDA block is plug-and-play and can be integrated into modern segmentation architectures beyond the UNet family.
\end{abstract}

\begin{IEEEkeywords}
medical image segmentation, dual attention, efficient transformers, low-rank attention, hardware-aware inference
\end{IEEEkeywords}


% ===========================================================================
\section{Introduction}
% ===========================================================================

Accurate medical image segmentation is a prerequisite for computer-aided diagnosis, treatment planning, and surgical guidance across a wide range of clinical scenarios including multi-organ CT analysis, skin lesion assessment, and endoscopic polyp detection~\cite{ronneberger2015unet}. Deep learning has driven successive performance improvements, evolving from fully convolutional networks to attention-augmented architectures. However, deploying high-capacity segmentation models at the point of care---on edge devices, mobile workstations, or low-memory GPUs---remains a practical bottleneck, since state-of-the-art models are commonly tuned for powerful cloud hardware and exhibit prohibitive inference latency on modest devices.

The U-Net family~\cite{ronneberger2015unet} established the encoder-decoder paradigm with skip connections and became the dominant framework for medical image segmentation. CNN-based successors such as U-Net++~\cite{zhou2018unetpp} and Attention U-Net~\cite{oktay2018attention} improved multi-scale feature reuse and attention-guided selection, but CNNs are inherently limited in modeling long-range spatial dependencies due to their local receptive fields. TransUNet~\cite{chen2021transunet} bridged this gap by embedding a Vision Transformer (ViT)~\cite{dosovitskiy2020image} encoder within a U-Net decoder, providing global context through self-attention over image patches. Swin-Unet~\cite{cao2022swinunet} extended this by replacing ViT with the hierarchical Swin Transformer~\cite{liu2021swin}, enabling multi-scale representation learning with shifted-window self-attention.

Despite these advances, Transformer-based methods still lack mechanisms to explicitly capture fine-grained spatial boundary details and inter-channel semantic correlations that are particularly important in medical imaging, where organ boundaries are subtle and channel responses encode tissue-type contrast. Dual Attention Network (DANet)~\cite{fu2019dual}, originally proposed for natural scene segmentation, addresses this gap through two complementary modules: a Position Attention Module (PAM) that computes pixel-wise spatial dependencies via global self-attention, and a Channel Attention Module (CAM) that captures cross-channel semantic correlations. DA-TransUNet~\cite{sun2024datransunet} is the first work to integrate dual attention blocks into both the encoder and skip connections of a Transformer U-Net, achieving state-of-the-art results on six medical segmentation benchmarks.

However, the dual attention mechanism carries two fundamental computational bottlenecks. For a feature map of spatial resolution $N = H \times W$ and channel dimension $C$, the PAM computes an $N \times N$ attention matrix at cost $\mathcal{O}(N^2 C)$, while the CAM computes a $C \times C$ matrix at cost $\mathcal{O}(C^2 N)$. For high-resolution feature maps or large channel counts, these quadratic terms become intractable. Furthermore, DA-TransUNet uses a fixed channel compression ratio of 1/16 for all configurations, which is tuned for high-end GPU training and offers no mechanism for adaptation to heterogeneous deployment environments.

Critically, while efficient attention mechanisms---including Swin-style window partitioning~\cite{liu2021swin} and low-rank matrix approximation as in Linformer~\cite{wang2020linformer}---have been extensively studied for vanilla Transformer self-attention, they have not been specifically applied to the dual attention block structure in a medical segmentation context. This gap motivates AdaDA-TransUNet.

Our contributions are summarized as follows:
\begin{enumerate}
    \item We propose a \textbf{Low-Rank Windowed PAM} that combines Swin-style window partitioning with low-rank key-value projection, reducing PAM complexity from $\mathcal{O}(N^2 C)$ to $\mathcal{O}(N r C)$ where $r \ll N$.
    \item We propose a \textbf{Grouped CAM} that partitions channels into $G$ independent groups before computing channel attention, reducing CAM complexity from $\mathcal{O}(C^2 N)$ to $\mathcal{O}(C^2 N / G)$.
    \item We introduce a \textbf{differentiable adaptive feature gate} that learns a per-channel soft blending ratio between PAM and CAM outputs, replacing the fixed static fusion in DA-TransUNet with a data-driven, gradient-trainable compression ratio.
    \item We design a \textbf{hardware-aware configuration module} that selects window size, rank, and group count at inference time based on available GPU memory, enabling graceful performance-efficiency trade-off without retraining.
\end{enumerate}


% ===========================================================================
\section{Related Work}
% ===========================================================================

\subsection{CNN-Based Medical Image Segmentation}

U-Net~\cite{ronneberger2015unet} introduced the encoder-decoder architecture with skip connections, which remains the foundational design for medical segmentation. U-Net++ \cite{zhou2018unetpp} proposed dense skip pathways with nested sub-networks to reduce the semantic gap between encoder and decoder features. Attention U-Net~\cite{oktay2018attention} introduced attention gates on skip connections, suppressing irrelevant background activations. While these architectures excel at local feature extraction, their CNN backbone restricts the receptive field and limits long-range contextual modeling.

\subsection{Transformer-Based Medical Segmentation}

TransUNet~\cite{chen2021transunet} was among the first to hybridize CNN feature extraction with ViT global attention, demonstrating that Transformer encoders capture complementary information to CNNs in medical images. Swin-Unet~\cite{cao2022swinunet} replaced the ViT encoder with a pure Swin Transformer encoder-decoder, achieving hierarchical feature representations with reduced complexity via shifted-window attention. UCTransNet~\cite{wang2022uctransnet} revisited skip connections through channel-wise cross-attention to address the semantic gap between encoder and decoder. These methods improve global context modeling but do not explicitly model the dual spatial-channel dependencies that are particularly beneficial in medical imaging.

\subsection{Efficient Attention Mechanisms}

The quadratic cost of self-attention has motivated extensive work on efficient alternatives. Swin Transformer~\cite{liu2021swin} introduced non-overlapping window partitioning to restrict attention computation to local windows, reducing spatial attention from $\mathcal{O}(N^2)$ to $\mathcal{O}(N M^2)$ where $M$ is the window size. Linformer~\cite{wang2020linformer} demonstrated that the self-attention matrix has a low stable rank, proposing to project keys and values from sequence length $N$ to rank $r$, reducing complexity to $\mathcal{O}(Nr)$. Performer~\cite{choromanski2020performer} approximates the softmax kernel with random features. These efficient attention methods have been applied almost exclusively to vanilla Transformer self-attention over patch tokens; their application to the dual attention block structure in DANet has not been explored.

\subsection{Dual Attention Networks}

DANet~\cite{fu2019dual} introduced dual attention for natural scene image segmentation, showing that jointly modeling spatial and channel dependencies leads to more discriminative feature representations. DAResUNet~\cite{shi2020daResunet} applied dual attention to ResUNet for medical image analysis. DA-DSUNet~\cite{tang2022daDSunet} extended dual attention to head-and-neck tumor segmentation. DA-TransUNet~\cite{sun2024datransunet} achieved the current state of the art by integrating DA-blocks into both the Transformer encoder and skip connections of a U-Net. None of these works addresses the efficiency limitations of the dual attention mechanism itself.


% ===========================================================================
\section{Proposed Method}
% ===========================================================================

\subsection{Architecture Overview}

AdaDA-TransUNet follows the overall encoder-decoder structure of DA-TransUNet, but replaces every standard DA-block with the proposed AdaDA block. The encoder is a hybrid CNN--Transformer backbone: an input image is first processed by three convolutional stages that produce multi-scale feature maps at $1/4$, $1/8$, and $1/16$ of the input resolution, followed by a ViT-B/16 Transformer that extracts global context at the $1/16$ scale. An AdaDA block is appended after each CNN encoder stage and after the Transformer stage. Skip connections from each encoder stage are passed through an AdaDA block before being merged with the decoder features, mirroring DA-TransUNet's use of dual attention to filter skip-connection redundancy. The decoder upsamples hierarchically from $1/16$ to $1/1$ resolution via transposed convolutions and concatenated skip features, with a final $1 \times 1$ convolution producing the segmentation mask. The architecture is depicted in Fig.~\ref{fig:architecture}.

\begin{figure}[t]
  \centering
  % \includegraphics[width=\columnwidth]{fig_architecture.pdf}
  \fbox{\rule{0pt}{3cm}\rule{\columnwidth}{0pt}}
  \caption{AdaDA-TransUNet architecture. AdaDA blocks replace all DA blocks in encoder stages, skip connections, and the bottleneck. The hardware-aware configuration module sets rank $r$, window size $M$, and group count $G$ based on available GPU memory.}
  \label{fig:architecture}
\end{figure}

\subsection{Low-Rank Windowed Position Attention Module}

\subsubsection{Complexity analysis of standard PAM}
Given an input feature map $\mathbf{A} \in \mathbb{R}^{C \times H \times W}$, the original PAM in DANet~\cite{fu2019dual} projects $\mathbf{A}$ to query $\mathbf{B}$, key $\mathbf{C}$, and value $\mathbf{D}$, each of shape $C \times N$ where $N = HW$. The spatial attention matrix $\mathbf{S} \in \mathbb{R}^{N \times N}$ is:
\begin{equation}
  S_{ji} = \frac{\exp\!\left(\mathbf{B}_i^\top \mathbf{C}_j\right)}{\sum_{k=1}^{N} \exp\!\left(\mathbf{B}_k^\top \mathbf{C}_j\right)},
  \label{eq:pam_attn}
\end{equation}
and the output is $E_j = \alpha \sum_{i=1}^{N} S_{ji} \mathbf{D}_i + \mathbf{A}_j$, where $\alpha$ is a learnable scalar. Computing $\mathbf{S}$ requires $\mathcal{O}(N^2 C)$ operations. For a $224 \times 224$ image at $1/4$ resolution, $N = 3136$, making the full attention matrix computationally prohibitive.

\subsubsection{Window partitioning}
Inspired by Swin Transformer~\cite{liu2021swin}, we partition the feature map into non-overlapping $M \times M$ windows. Let $n_W = (H/M)(W/M)$ be the number of windows. Each window contains $N_w = M^2$ tokens. Applying PAM independently within each window reduces complexity from $\mathcal{O}(N^2 C)$ to $\mathcal{O}(n_W \cdot N_w^2 \cdot C) = \mathcal{O}(N M^2 C)$. Since $M \ll \sqrt{N}$, this is a substantial reduction, though it sacrifices global spatial dependencies beyond window boundaries.

\subsubsection{Low-rank key-value projection}
Following Linformer~\cite{wang2020linformer}, we observe that the within-window attention matrix has low intrinsic rank due to spatial redundancy. We project keys and values from $N_w = M^2$ to a rank $r \ll M^2$ using a learned linear projection $\mathbf{W}_r \in \mathbb{R}^{r \times N_w}$:
\begin{align}
  \tilde{\mathbf{C}} &= \mathbf{W}_r \mathbf{C}^\top \in \mathbb{R}^{r \times C}, \label{eq:low_rank_key}\\
  \tilde{\mathbf{D}} &= \mathbf{W}_r \mathbf{D}^\top \in \mathbb{R}^{r \times C}. \label{eq:low_rank_val}
\end{align}
The attention scores become $\mathbf{B}^\top \tilde{\mathbf{C}}^\top \in \mathbb{R}^{N_w \times r}$, reducing the within-window cost from $\mathcal{O}(N_w^2 C)$ to $\mathcal{O}(N_w r C)$. The total complexity across all windows is then $\mathcal{O}(N r C)$, where $r \ll N_w \ll N$. This corresponds to the low-rank approximation $\mathbf{A} \approx \mathbf{P}\mathbf{Q}$ where $\mathbf{P} \in \mathbb{R}^{N \times r}$ and $\mathbf{Q} \in \mathbb{R}^{r \times N}$.

The complete LowRankWindowedPAM forward pass is formalized in Algorithm~\ref{alg:lrwpam}.

\begin{algorithm}[t]
\caption{LowRankWindowedPAM Forward Pass}
\label{alg:lrwpam}
\begin{algorithmic}[1]
\REQUIRE Feature map $\mathbf{x} \in \mathbb{R}^{B \times C \times H \times W}$, window size $M$, rank $r$
\STATE $\mathbf{x}_w \leftarrow \texttt{window\_partition}(\mathbf{x}, M)$
  \hfill \textit{\small // $(B n_W, C, M^2)$}
\STATE $\mathbf{B}, \mathbf{C}, \mathbf{D} \leftarrow \text{Conv1d projections of } \mathbf{x}_w$
  \hfill \textit{\small // $(B n_W, C, N_w)$}
\STATE $\tilde{\mathbf{C}} \leftarrow \mathbf{W}_r(\mathbf{C})$,\quad $\tilde{\mathbf{D}} \leftarrow \mathbf{W}_r(\mathbf{D})$
  \hfill \textit{\small // $(B n_W, C, r)$}
\STATE $\mathbf{S} \leftarrow \text{softmax}\!\left(\mathbf{B}^\top \tilde{\mathbf{C}}\right)$
  \hfill \textit{\small // $(B n_W, N_w, r)$}
\STATE $\mathbf{E} \leftarrow \mathbf{S} \tilde{\mathbf{D}}^\top + \alpha\, \mathbf{x}_w$
  \hfill \textit{\small // $(B n_W, C, N_w)$}
\STATE \textbf{return} $\texttt{window\_reverse}(\mathbf{E}, M, H, W)$
\end{algorithmic}
\end{algorithm}


\subsection{Grouped Channel Attention Module}

\subsubsection{Standard CAM}
The original CAM~\cite{fu2019dual} reshapes $\mathbf{A}$ to $\mathbb{R}^{C \times N}$ and computes:
\begin{equation}
  X_{ji} = \frac{\exp(\mathbf{A}_i \cdot \mathbf{A}_j)}{\sum_{k=1}^{C} \exp(\mathbf{A}_k \cdot \mathbf{A}_j)},
  \label{eq:cam_attn}
\end{equation}
with output $E_j = \beta \sum_{i=1}^{C} X_{ji} \mathbf{A}_i + \mathbf{A}_j$, where $\beta$ is a learnable scalar. Computing the $C \times C$ matrix costs $\mathcal{O}(C^2 N)$.

\subsubsection{Grouped decomposition}
We partition the $C$ channels into $G$ equal-sized groups, each of width $C_g = C / G$. Within each group $g$, we compute an independent per-group channel attention matrix $\mathbf{X}^{(g)} \in \mathbb{R}^{C_g \times C_g}$:
\begin{equation}
  X^{(g)}_{ji} = \frac{\exp\!\left(\mathbf{A}^{(g)}_i \cdot \mathbf{A}^{(g)}_j\right)}{\sum_{k=1}^{C_g} \exp\!\left(\mathbf{A}^{(g)}_k \cdot \mathbf{A}^{(g)}_j\right)},
  \label{eq:grouped_cam}
\end{equation}
where $\mathbf{A}^{(g)} \in \mathbb{R}^{C_g \times N}$ denotes the feature slice for group $g$. The total cost across $G$ groups is $\mathcal{O}(G \cdot C_g^2 N) = \mathcal{O}(C^2 N / G)$. The output channels are concatenated to restore the original channel dimension. Grouped channel attention preserves the semantics of the original CAM (each group computes softmax attention over its own channel set) while reducing complexity by a factor $G$. This is well-suited to multi-class medical segmentation where different channel subsets may encode distinct tissue or pathology categories.

\subsection{Adaptive Feature Gate (Learnable Compression Ratio)}

DA-TransUNet fuses PAM and CAM outputs by summation after a fixed 1/16 channel compression, treating position and channel cues as equally important for all inputs and all spatial locations. We instead propose a differentiable per-channel gate $\mathbf{g} \in (0, 1)^C$ that adapts the fusion ratio based on global context:
\begin{equation}
  \mathbf{g} = \sigma\!\left(\mathbf{W}_{\text{gate}}\, \text{AvgPool}(\mathbf{x})\right) \in (0, 1)^C,
  \label{eq:gate}
\end{equation}
where $\mathbf{W}_{\text{gate}} \in \mathbb{R}^{C \times C}$ is a learned linear layer and $\sigma$ is the sigmoid function. The fused feature is:
\begin{equation}
  \mathbf{y} = \mathbf{W}_f \!\left(\mathbf{g} \odot \hat{\mathbf{P}} + (1 - \mathbf{g}) \odot \hat{\mathbf{C}}\right) + \mathbf{x},
  \label{eq:fusion}
\end{equation}
where $\hat{\mathbf{P}}$ and $\hat{\mathbf{C}}$ are the PAM and CAM outputs, $\odot$ is broadcast element-wise multiplication, and $\mathbf{W}_f$ is a $1 \times 1$ convolution. The gate is fully differentiable—no discrete indexing or \texttt{.item()} detachment—so gradients flow through both the gate and the attended features. When $g_c \to 1$, channel $c$ relies on spatial context from PAM; when $g_c \to 0$, channel $c$ relies on cross-channel context from CAM. This allows the model to automatically discover which feature types are most informative per channel and per input, functioning as a learned compression ratio that subsumes the static setting in DA-TransUNet.

\subsection{Hardware-Aware Configuration Module}

For deployment on diverse hardware from high-memory data center GPUs to resource-constrained edge devices, we introduce a hardware-aware configuration module that selects hyperparameters $\{r, M, G\}$ at inference time based on available GPU memory:

\begin{equation}
  \{r, M, G\} =
  \begin{cases}
    \{64,\;14,\;4\} & \text{if } F > 8\text{ GB}, \\
    \{32,\;\;7,\;8\} & \text{if } 4 < F \leq 8\text{ GB}, \\
    \{16,\;\;7,\;16\} & \text{otherwise},
  \end{cases}
  \label{eq:hw_config}
\end{equation}
where $F$ is free GPU memory in GB. Higher available memory permits larger rank (more expressive low-rank approximation), larger windows (wider spatial context), and fewer groups (richer channel interactions). Lower available memory activates a leaner configuration that maintains functionality while fitting within device constraints. Crucially, the model is trained once with a default configuration and the configuration module only adjusts the three structural hyperparameters at inference; no retraining or fine-tuning is required for deployment on a new device class.

\subsection{Loss Function}

We adopt the combined Dice--Cross-Entropy loss with equal weighting, consistent with DA-TransUNet:
\begin{equation}
  \mathcal{L} = \frac{1}{2}\mathcal{L}_{\text{Dice}} + \frac{1}{2}\mathcal{L}_{\text{CE}},
  \label{eq:loss}
\end{equation}
where $\mathcal{L}_{\text{Dice}} = 1 - \frac{2 \sum_i p_i y_i}{\sum_i p_i + \sum_i y_i}$ promotes overlap between predicted and ground-truth masks, and $\mathcal{L}_{\text{CE}}$ provides per-pixel class probability supervision.

\subsection{Complexity Summary}

Table~\ref{tab:complexity} summarizes the theoretical complexity of each attention module.

\begin{table}[t]
  \caption{Complexity Comparison of Attention Modules}
  \label{tab:complexity}
  \begin{center}
  \begin{tabular}{lcc}
    \toprule
    Module & Original & Proposed \\
    \midrule
    PAM & $\mathcal{O}(N^2 C)$ & $\mathcal{O}(N r C)$ \\
    CAM & $\mathcal{O}(C^2 N)$ & $\mathcal{O}(C^2 N / G)$ \\
    Fusion & Static & Learnable gate \\
    Config & Fixed & Hardware-aware \\
    \bottomrule
  \end{tabular}
  \end{center}
  \vspace{-0.5em}
  \footnotesize{$N$: spatial tokens, $C$: channels, $r$: rank ($r \ll N$), $G$: groups.}
\end{table}


% ===========================================================================
\section{Experiments}
% ===========================================================================

\subsection{Datasets}

We evaluate AdaDA-TransUNet on three benchmark medical image segmentation datasets spanning CT, dermatology, and endoscopy modalities.

\textbf{Synapse Multi-Organ CT Segmentation.} The Synapse dataset~\cite{landman2015miccai} contains 30 abdominal CT scans with ground-truth annotations for 8 organs (aorta, gallbladder, spleen, left/right kidney, liver, stomach, pancreas). Following the standard split used by TransUNet and DA-TransUNet, we use 18 scans for training and 12 for testing. Evaluation metrics are mean Dice Score Coefficient (DSC) and 95th percentile Hausdorff Distance (HD95).

\textbf{ISIC 2018 Skin Lesion Segmentation.} The ISIC 2018 challenge dataset~\cite{codella2019isic} provides 2594 dermoscopic images with pixel-level lesion annotations. We use the standard split of 1815 training, 259 validation, and 520 test images. We report DSC and mean intersection-over-union (mIoU).

\textbf{Kvasir-SEG Polyp Segmentation.} Kvasir-SEG~\cite{jha2020kvasir} contains 1000 endoscopy images with polyp annotations. We use an 800/200 train/test split and report DSC and mIoU.

\subsection{Implementation Details}

All experiments use a hybrid CNN--ViT-B/16 encoder pretrained on ImageNet-21K. The CNN encoder consists of three ResNet-like stages producing features at strides 4, 8, and 16. The Transformer has 12 layers with hidden dimension 768 and 12 attention heads. AdaDA blocks use default hyperparameters $r=32$, $M=7$, $G=8$ during training (corresponding to the mid-tier hardware configuration in \eqref{eq:hw_config}).

Following DA-TransUNet, we train with SGD (momentum 0.9, weight decay $10^{-4}$) with initial learning rate 0.01 and polynomial decay schedule. Batch size is 24 and input images are resized to $224 \times 224$ with random horizontal flipping and random rotation for data augmentation. Training runs for 150 epochs. All experiments are conducted on a single NVIDIA Tesla T4 (15 GB VRAM). We report wall-clock training time and peak VRAM to compare hardware efficiency alongside segmentation accuracy.

To validate our experimental setup, we reproduced DA-TransUNet under identical conditions (same hardware, optimizer, and random seed). Our reproduction achieves 79.36\% DSC and 26.64 mm HD95, within 2\% DSC of the reported 81.03\%, confirming implementation correctness. The small gap is attributable to the absence of validation-based checkpoint selection in our setup; we use the final epoch checkpoint (epoch 149) for both models to ensure a fair comparison.

\subsection{Comparison with State-of-the-Art}

\subsubsection{Synapse multi-organ CT}
Table~\ref{tab:synapse} reports DSC and HD95 on the Synapse dataset. AdaDA-TransUNet achieves competitive Dice and HD95 compared to DA-TransUNet while significantly reducing FLOPs (see Table~\ref{tab:efficiency}).

\begin{table}[t]
  \caption{Segmentation Results on Synapse Multi-Organ CT}
  \label{tab:synapse}
  \begin{center}
  \setlength{\tabcolsep}{4pt}
  \begin{tabular}{lcc}
    \toprule
    Method & DSC (\%) $\uparrow$ & HD95 (mm) $\downarrow$ \\
    \midrule
    U-Net~\cite{ronneberger2015unet}         & 74.68 & 36.87 \\
    Attention U-Net~\cite{oktay2018attention} & 77.77 & 36.02 \\
    TransUNet~\cite{chen2021transunet}        & 77.48 & 31.69 \\
    Swin-Unet~\cite{cao2022swinunet}          & 79.13 & 21.55 \\
    UCTransNet~\cite{wang2022uctransnet}      & 78.23 & 26.75 \\
    DA-TransUNet~\cite{sun2024datransunet}    & 81.03 & 17.84 \\
    DA-TransUNet$^\dagger$ (reproduced)       & 79.36 & 26.64 \\
    \midrule
    \textbf{AdaDA-TransUNet (ours)}           & \textbf{--}   & \textbf{--}   \\
    \bottomrule
  \end{tabular}
  \end{center}
  \vspace{-0.4em}
  \footnotesize{Baselines from published papers. DA-TransUNet† denotes our reproduction (79.36\% vs.\ published 81.03\%); the gap is attributable to absence of validation-based checkpoint selection---both our models use the final-epoch checkpoint for a fair comparison. -- AdaDA result pending.}
\end{table}

\subsubsection{Skin lesion and polyp datasets}
Table~\ref{tab:other_datasets} compares DSC and mIoU on ISIC 2018 and Kvasir-SEG.

\begin{table}[t]
  \caption{Segmentation Results on ISIC 2018 and Kvasir-SEG}
  \label{tab:other_datasets}
  \begin{center}
  \setlength{\tabcolsep}{4pt}
  \begin{tabular}{lcccc}
    \toprule
    \multirow{2}{*}{Method} & \multicolumn{2}{c}{ISIC 2018} & \multicolumn{2}{c}{Kvasir-SEG} \\
    \cmidrule(lr){2-3}\cmidrule(lr){4-5}
    & DSC (\%) & IoU (\%) & DSC (\%) & IoU (\%) \\
    \midrule
    U-Net~\cite{ronneberger2015unet}            & 86.7 & 76.8 & 81.8 & 74.6 \\
    TransUNet~\cite{chen2021transunet}           & 88.4 & 79.8 & 84.7 & 78.3 \\
    DA-TransUNet~\cite{sun2024datransunet}       & 90.3 & 82.5 & 87.6 & 80.5 \\
    \midrule
    \textbf{AdaDA-TransUNet (ours)}             & --   & --   & --   & -- \\
    \bottomrule
  \end{tabular}
  \end{center}
  \vspace{-0.4em}
  \footnotesize{Baseline values from DA-TransUNet~\cite{sun2024datransunet}. -- Pending experiment results.}
\end{table}

\subsubsection{Efficiency comparison}
Table~\ref{tab:efficiency} compares model parameters, GFLOPs, training VRAM, training time, and per-volume inference time to highlight the efficiency gains of AdaDA-TransUNet.

\begin{table*}[t]
  \caption{Efficiency Comparison on Synapse, NVIDIA Tesla T4 (batch=24, 150 epochs)}
  \label{tab:efficiency}
  \begin{center}
  \setlength{\tabcolsep}{5pt}
  \begin{tabular}{lccccccc}
    \toprule
    Method & Params (M) & GFLOPs & Train Time & Train VRAM & Infer Time (s/vol) & Infer VRAM & Multi-GPU \\
    \midrule
    TransUNet~\cite{chen2021transunet}     & 105.3       & 57.4        & --               & --               & --               & --               & No \\
    Swin-Unet~\cite{cao2022swinunet}       & 41.4        & 11.5        & --               & --               & --               & --               & No \\
    DA-TransUNet~\cite{sun2024datransunet} & 107.95      & \textit{XX} & 6.3h (T4$\times$1) & \textit{XX}   & 112.4            & 0.5 GB           & No$^\ddagger$ \\
    \midrule
    \multirow{2}{*}{\textbf{AdaDA ($r$=32, $M$=7, $G$=8)}}
      & \multirow{2}{*}{\textit{XX}} & \multirow{2}{*}{\textit{XX}}
      & \textit{XX}h (T4$\times$1) & \textit{XX} & \multirow{2}{*}{\textit{XX}} & \multirow{2}{*}{\textit{XX}} & No \\
      & & & \textit{XX}h (T4$\times$2)$^\S$ & \textit{XX} & & & Yes \\
    \bottomrule
  \end{tabular}
  \end{center}
  \vspace{-0.4em}
  \footnotesize{GFLOPs by \texttt{thop} on a single $224\!\times\!224$ slice. Train VRAM = peak during backward pass; Infer Time = avg seconds per CT volume over 12 test cases (includes NIfTI I/O); Infer VRAM = peak during forward pass. TransUNet/Swin-Unet GFLOPs from published papers. \textit{XX} = pending.\\
  $^\ddagger$~DA-TransUNet OOMs on 2$\times$T4 DataParallel: the $\mathcal{O}(N^2)$ PAM backward stores a $3136\!\times\!3136$ attention matrix per sample at the 56$\times$56 skip connection, exhausting GPU\,0 memory when it aggregates gradients from both devices.\\
  $^\S$~AdaDA's windowed low-rank PAM reduces the per-sample attention footprint by $>$25$\times$ at the same resolution, enabling 2-GPU DataParallel training and approximately halving wall-clock time with no change to model weights or accuracy.}
\end{table*}

Additionally, DA-TransUNet cannot utilize multi-GPU data-parallel training due to out-of-memory errors during backward pass of the $\mathcal{O}(N^2)$ PAM attention matrix (even at batch size 16 per GPU on T4). In contrast, AdaDA-TransUNet's windowed low-rank attention reduces the per-GPU memory footprint sufficiently to enable 2-GPU DataParallel training, approximately halving wall-clock training time on equivalent hardware.

\subsection{Ablation Study}

To quantify the contribution of the key design choices, we perform targeted ablation experiments on the Synapse dataset. Table~\ref{tab:ablation} examines the effect of the adaptive gate and the low-rank projection dimension $r$.

\begin{table*}[t]
  \caption{Ablation Study on Synapse Multi-Organ CT (T4$\times$2, SGD, lr=0.01, batch=24, 150 epochs, seed=1234)}
  \label{tab:ablation}
  \begin{center}
  \setlength{\tabcolsep}{5pt}
  \begin{tabular}{lcccccc}
    \toprule
    Configuration & DSC (\%) $\uparrow$ & HD95 (mm) $\downarrow$ & GFLOPs & Train VRAM & Infer Time (s/vol) & Infer VRAM \\
    \midrule
    DA-TransUNet (baseline)    & 79.36        & 26.64        & \textit{XX} & \textit{XX} & 112.4        & 0.5 GB       \\
    AdaDA, fixed gate ($r$=32) & \textit{XX}  & \textit{XX}  & \textit{XX} & \textit{XX} & \textit{XX} & \textit{XX} \\
    AdaDA, ada-gate ($r$=8)    & \textit{XX}  & \textit{XX}  & \textit{XX} & \textit{XX} & \textit{XX} & \textit{XX} \\
    \textbf{AdaDA, ada-gate ($r$=32) [ours]} & \textit{XX} & \textit{XX} & \textit{XX} & \textit{XX} & \textit{XX} & \textit{XX} \\
    \bottomrule
  \end{tabular}
  \end{center}
  \vspace{-0.4em}
  \footnotesize{\textbf{All rows trained under identical conditions} for a controlled comparison. The DA-TransUNet baseline uses our reproduced result (79.36\%) rather than the published 81.03\%, which used different hardware and validation-based checkpoint selection; using our reproduction ensures $\Delta$DSC between rows reflects only architectural differences. GFLOPs by \texttt{thop} on a single $224\!\times\!224$ slice. \textit{XX} = pending.}
\end{table*}

Comparing rows 1 and 4 isolates the combined effect of all AdaDA components. Rows 2 vs.\ 4 quantify the contribution of the adaptive gate: replacing the learned gate with a fixed 0.5 blend shows how much the gradient-trainable fusion ratio contributes beyond the windowed and grouped attention alone. Rows 3 vs.\ 4 show the sensitivity to rank: $r$=8 is more memory-efficient but may sacrifice expressiveness, while $r$=32 is our default mid-tier hardware configuration.


% ===========================================================================
\section{Conclusion}
% ===========================================================================

We presented AdaDA-TransUNet, a hardware-aware adaptive dual attention Transformer UNet for efficient medical image segmentation. By targeting efficiency improvements directly at the dual attention block---through low-rank windowed position attention, grouped channel attention, a differentiable adaptive feature gate, and a hardware-aware configuration module---our method addresses a gap not covered by prior efficient Transformer work, which has focused predominantly on vanilla self-attention over patch tokens.

The proposed modules reduce PAM complexity from $\mathcal{O}(N^2 C)$ to $\mathcal{O}(NrC)$ and CAM complexity from $\mathcal{O}(C^2 N)$ to $\mathcal{O}(C^2 N / G)$, while the adaptive gate preserves the expressiveness of dual attention by learning a per-channel fusion ratio from data. The hardware-aware configuration enables deployment across a wide range of GPU memory budgets without retraining.

As future work, we plan to extend AdaDA-TransUNet to (1) multi-modal segmentation where cross-modality DA-blocks can align spatial and channel cues across MRI and CT; and (2) video-based lesion tracking where temporal continuity across frames can be incorporated into the position attention window, eliminating redundant per-frame computation.

\section*{Acknowledgment}

The authors thank the open-source contributors of DA-TransUNet and Swin Transformer for their publicly available codebases.


% ===========================================================================
\begin{thebibliography}{00}

\bibitem{ronneberger2015unet}
O. Ronneberger, P. Fischer, and T. Brox, ``U-Net: Convolutional networks for biomedical image segmentation,'' in \textit{Proc. MICCAI}, 2015, pp. 234--241.

\bibitem{zhou2018unetpp}
Z. Zhou, M. M. R. Siddiquee, N. Tajbakhsh, and J. Liang, ``UNet++: A nested U-Net architecture for medical image segmentation,'' in \textit{DLMIA/ML-CDS @ MICCAI}, 2018, pp. 3--11.

\bibitem{oktay2018attention}
O. Oktay \textit{et al.}, ``Attention U-Net: Learning where to look for the pancreas,'' in \textit{Proc. MIDL}, 2018.

\bibitem{chen2021transunet}
J. Chen \textit{et al.}, ``TransUNet: Transformers make strong encoders for medical image segmentation,'' \textit{arXiv:2102.04306}, 2021.

\bibitem{dosovitskiy2020image}
A. Dosovitskiy \textit{et al.}, ``An image is worth 16x16 words: Transformers for image recognition at scale,'' in \textit{Proc. ICLR}, 2021.

\bibitem{liu2021swin}
Z. Liu \textit{et al.}, ``Swin Transformer: Hierarchical vision transformer using shifted windows,'' in \textit{Proc. ICCV}, 2021, pp. 10012--10022.

\bibitem{cao2022swinunet}
H. Cao \textit{et al.}, ``Swin-Unet: Unet-like pure transformer for medical image segmentation,'' in \textit{Proc. ECCV Workshops}, 2022.

\bibitem{wang2022uctransnet}
H. Wang \textit{et al.}, ``UCTransNet: Rethinking the skip connections in U-Net from a channel-wise perspective with transformer,'' in \textit{Proc. AAAI}, 2022, pp. 2441--2449.

\bibitem{fu2019dual}
J. Fu \textit{et al.}, ``Dual attention network for scene segmentation,'' in \textit{Proc. CVPR}, 2019, pp. 3146--3154.

\bibitem{sun2024datransunet}
G. Sun \textit{et al.}, ``DA-TransUNet: Integrating spatial and channel dual attention with transformer U-Net for medical image segmentation,'' \textit{Front. Bioeng. Biotechnol.}, vol. 12, p. 1398237, 2024.

\bibitem{wang2020linformer}
S. Wang \textit{et al.}, ``Linformer: Self-attention with linear complexity,'' \textit{arXiv:2006.04768}, 2020.

\bibitem{choromanski2020performer}
K. Choromanski \textit{et al.}, ``Rethinking attention with Performers,'' in \textit{Proc. ICLR}, 2021.

\bibitem{shi2020daResunet}
W. Shi, J. Li, and J. Wu, ``DAResUNet: Dual attention residual UNet for retinal vessel segmentation,'' in \textit{Proc. BIBM}, 2020.

\bibitem{tang2022daDSunet}
H. Tang \textit{et al.}, ``DA-DSUNet: Dual attention-based deep supervision UNet for head-and-neck tumor segmentation,'' \textit{Neurocomputing}, vol. 515, pp. 200--211, 2022.

\bibitem{landman2015miccai}
B. Landman \textit{et al.}, ``MICCAI multi-atlas labeling beyond the cranial vault workshop,'' in \textit{MICCAI Workshop}, 2015.

\bibitem{codella2019isic}
N. Codella \textit{et al.}, ``Skin lesion analysis toward melanoma detection 2018: A challenge hosted by the International Skin Imaging Collaboration,'' \textit{arXiv:1902.03368}, 2019.

\bibitem{shiraishi2000jsrt}
J. Shiraishi \textit{et al.}, ``Development of a digital image database for chest radiographs with and without a lung nodule,'' \textit{Am. J. Roentgenol.}, vol. 174, no. 1, pp. 71--74, 2000.

\bibitem{bernal2015cvcclinicdb}
J. Bernal, F. J. Sanchez, G. Fernandez-Esparrach, D. Gil, C. Rodriguez, and F. Vilarino, ``WM-DOVA maps for accurate polyp highlighting in colonoscopy,'' \textit{Comput. Med. Imaging Graph.}, vol. 43, pp. 99--111, 2015.

\bibitem{jha2020kvasir}
D. Jha \textit{et al.}, ``Kvasir-SEG: A segmentation dataset and benchmark for gastrointestinal polyp segmentation,'' in \textit{Proc. MMM}, 2020, pp. 451--462.

\end{thebibliography}

\end{document}
